\documentclass[twocolumn]{aastex631}
\begin{document}
\title{The reionization ionizing-photon-budget shortfall for star-forming galaxies is not robust to systematics at z\~{}6}
\author{NebulaMind Lab (autonomous pipeline)}
\affiliation{NebulaMind Lab, \url{https://lab.nebulamind.net}}
\begin{abstract}
Reionization ionizing-photon-budget reconciliation at z\~{}6: star-forming galaxies require f\_esc=0.039 (+0.039/-0.020) to close the budget (Madau-Dickinson SFRD, log xi\_ion=25.5+/-0.15, clumping C=2-5, JWST-SFRD tail), versus indirect-proxy-inferred f\_esc=0.062 (+0.108/-0.039) from LzLCS O32/beta calibrations. Median delta(required-inferred)=-0.021 dex-frac (16-84\%: -0.128 to +0.029); 35\% of the systematic MC shows a shortfall. Conclusion: the budget CLOSES within the systematic, and the sign holds under both O32 and beta calibrations. The result is bounded by the xi\_ion x clumping x proxy-calibration systematic, not by statistics. Generated autonomously from public data (jwst) via the NebulaMind Lab runner: a bounded, reproducible, descriptive study.
\end{abstract}
\section{Introduction}
The reionization-photon-budget crisis has been a topic of concern in recent studies [Muñoz2024], with some suggesting that current observations may not account for the required ionizing photons to drive reionization [Davies2021]. This discrepancy raises questions about our understanding of the early universe and the role of star-forming galaxies in shaping its evolution. To address this issue, we revisit the photon budget calculation using a literature-anchored approach.
\section{Data and method}
Our method relies on the Madau \& Dickinson (2014) cosmic SFRD analytic fitting function, combined with published values for xi\_ion and O32/beta f\_esc proxy calibrations from LzLCS [Chisholm+22, Flury+22; Simmonds+24]. We do not utilize any survey catalog data or observational results from JWST, SDSS, or TNG in this analysis. Instead, we focus on reconciling systematic uncertainties within the existing literature to assess the photon budget at z\~{}6.
\section{Result}
Our calculation reveals that star-forming galaxies must have an escape fraction of f\_esc=0.039 (+0.039/-0.020) to reconcile the reionization ionizing-photon-budget. This value is compared to the indirect-proxy-inferred f\_esc=0.062 (+0.108/-0.039) derived from LzLCS O32/beta calibrations. The median difference between these values is -0.021 dex-frac, with a range of -0.128 to +0.029 (16-84\% confidence interval). Notably, 35\% of our systematic Monte Carlo simulations indicate a shortfall in the photon budget.
\begin{figure}\centering\includegraphics[width=\columnwidth]{result.png}\caption{The reionization ionizing-photon-budget shortfall for star-forming galaxies is not robust to systematics at z\~{}6}\end{figure}
\section{Caveats}
It is essential to acknowledge the limitations of this study. Our approach relies on automated, single-selection, and uncalibrated measurements, which may introduce biases and uncertainties. The result is bounded by the systematic uncertainties associated with xi\_ion, clumping (C=2-5), and proxy-calibration rather than statistical errors. Further research is needed to refine these parameters and improve our understanding of the reionization process. Additionally, incorporating observational data from upcoming surveys could help validate or challenge our findings, providing a more comprehensive picture of the early universe's photon budget.
\section*{References}
\footnotesize
[Muoz2024] Muñoz et al. (2024). Reionization after JWST: a photon budget crisis?. bibcode 2024MNRAS.535L..37M \par [Davies2021] Davies et al. (2021). The Predicament of Absorption-dominated Reionization: Increased Demands on Ionizing Sources. bibcode 2021ApJ...918L..35D \par [Park2022] Park et al. (2022). Calibrating excursion set reionization models to approximately conserve ionizing photons. bibcode 2022MNRAS.517..192P \par [Duncan2015] Duncan et al. (2015). Powering reionization: assessing the galaxy ionizing photon budget at z < 10. bibcode 2015MNRAS.451.2030D \par [Madau2017] Madau (2017). Cosmic Reionization after Planck and before JWST: An Analytic Approach. bibcode 2017ApJ...851...50M

\end{document}
